```latex
\documentclass{article}
\usepackage{amsmath,amssymb,amsthm}
\usepackage{hyperref}

\newtheorem{theorem}{Theorem}
\newtheorem{lemma}{Lemma}
\newtheorem{remark}{Remark}

\theoremstyle{definition}
\newtheorem{definition}{Definition}

\newcommand{\cD}{\mathcal{D}}
\newcommand{\cC}{\mathcal{C}}
\newcommand{\cL}{\mathcal{L}}
\newcommand{\R}{\mathbb{R}}

\title{A regularity condition forces zero topological entropy}
\author{}
\date{}

\begin{document}

\maketitle

\begin{abstract}
Let $M$ be a compact connected smooth manifold, and let $X$ be a smooth
vector field on $M$. We prove that if every distributional solution $u$ of
$\cL_X u=f$ with $f\in C^\infty(M)$ is necessarily smooth, then the flow
generated by $X$ has zero topological entropy.
\end{abstract}

\section{Introduction and statement}

Let $M$ be a compact connected smooth manifold, and let $X\in
\mathfrak{X}^\infty(M)$ be a smooth vector field. Denote by
$\varphi_t:M\to M$ the flow generated by $X$, so that
\[
\left.\frac{d}{dt}\right|_{t=0}\varphi_t(x)=X(x).
\]
The Lie derivative of a distribution $u\in \cD'(M)$ along $X$ is defined by
duality:
\[
\langle \cL_X u,\psi\rangle
=
-\langle u,X\psi\rangle,
\qquad
\psi\in C^\infty(M).
\]
We assume the following regularity property:
\begin{quote}
If $f\in C^\infty(M)$ and $u\in \cD'(M)$ satisfy
$\cL_X u=f$, then $u\in C^\infty(M)$.
\end{quote}

The aim of this note is to prove the following result.

\begin{theorem}\label{thm:main}
Let $M$ be a compact connected smooth manifold and let $X$ be a smooth
vector field on $M$. If the above regularity property holds for $X$, then the
flow $\varphi_t$ generated by $X$ has zero topological entropy:
\[
h_{\mathrm{top}}(\varphi_1)=0.
\]
\end{theorem}

Here, as usual, the topological entropy of the flow is by definition the
topological entropy of its time-one map $\varphi_1$.

\section{Two elementary consequences of the regularity property}

\begin{lemma}\label{lem:no-singularities}
If the regularity property holds, then $X$ has no zeros.
\end{lemma}

\begin{proof}
Assume that $X(p)=0$ for some $p\in M$. Let $u=\delta_p\in\cD'(M)$ be the
Dirac mass at $p$. For every test function $\psi\in C^\infty(M)$,
\[
\langle \cL_X \delta_p,\psi\rangle
=
-\langle \delta_p,X\psi\rangle
=
-(X\psi)(p)
=
0,
\]
because $X(p)=0$. Thus $\cL_X u=0\in C^\infty(M)$. By the regularity
property, $u$ must be a smooth function. But $\delta_p$ is not a smooth
function when $\dim M\ge 1$. If $\dim M=0$, the statement is trivial, so we
assume $\dim M\ge 1$. This contradiction proves the lemma.
\end{proof}

\begin{lemma}\label{lem:invariant-distributions}
If the regularity property holds, then every invariant distribution of the
flow is smooth. More precisely, if $u\in\cD'(M)$ satisfies
$\cL_X u=0$, then $u\in C^\infty(M)$.
\end{lemma}

\begin{proof}
Apply the regularity property with $f=0$.
\end{proof}

Thus, to prove Theorem~\ref{thm:main}, it suffices to show that positive
topological entropy forces the existence of a non-smooth invariant
distribution.

\section{Entropy, Lyapunov exponents, and invariant distributions}

We shall use the following standard facts from smooth ergodic theory.

\begin{lemma}[Variational principle for flows]\label{lem:variational}
If $h_{\mathrm{top}}(\varphi_1)>0$, then there exists an ergodic
$\varphi_t$-invariant Borel probability measure $\mu$ on $M$ such that
\[
h_\mu(\varphi_1)>0.
\]
\end{lemma}

\begin{proof}
This is the usual variational principle for the homeomorphism $\varphi_1$,
together with the standard averaging argument over the flow. A reference is
\cite{KH}.
\end{proof}

\begin{lemma}[Ruelle inequality]\label{lem:ruelle}
Let $\mu$ be a $\varphi_t$-invariant ergodic probability measure, and let
$\lambda_1,\dots,\lambda_n$ be the Lyapunov exponents of $\varphi_1$ with
respect to $\mu$. Then
\[
h_\mu(\varphi_1)
\le
\int_M \sum_{\lambda_i(x)>0} \lambda_i(x)\,d\mu(x).
\]
In particular, if $h_\mu(\varphi_1)>0$, then $\mu$ has at least one positive
Lyapunov exponent.
\end{lemma}

\begin{proof}
This is Ruelle's inequality for $C^1$ maps; see \cite{KH, Rue}.
\end{proof}

The decisive dynamical input is the following invariant-distribution theorem.

\begin{lemma}[Anosov--Pesin invariant distribution theorem]\label{lem:anosov-pesin}
Let $\varphi_t$ be a $C^\infty$ flow on a compact manifold $M$, and let
$\mu$ be an ergodic $\varphi_t$-invariant probability measure. If $\mu$ has
a positive Lyapunov exponent, then there exists a non-smooth invariant
distribution
\[
u\in \cD'(M)\setminus C^\infty(M)
\]
such that $\cL_X u=0$.
\end{lemma}

\begin{proof}[Proof sketch and references]
By the Pesin unstable manifold theorem \cite{Pesin}, a positive Lyapunov
exponent gives a non-trivial unstable lamination for $\mu$. Ruelle and
Sullivan \cite{RS} constructed from such a lamination a non-zero invariant
current $C$ of positive degree, namely the unstable current. Since the
unstable direction is transverse to the flow direction $X$, contracting
$C$ repeatedly with $X$ produces a non-zero invariant distribution
$u=\iota_X^k C$, where $k$ is the dimension of the unstable bundle. Because
$C$ is not a smooth current, the resulting distribution $u$ is not smooth.
By construction,
\[
\cL_X u=\cL_X(\iota_X^k C)=\iota_X^k \cL_X C=0.
\]
Thus $u$ is the required non-smooth invariant distribution. Full details of
the Ruelle--Sullivan current construction can be found in \cite{RS}, and
the nonuniform hyperbolic version follows from Pesin theory \cite{Pesin}.
\end{proof}

\begin{remark}
The use of invariant distributions rather than periodic orbit measures is
essential: there exist minimal smooth flows with positive topological
entropy, hence with no periodic orbits at all. For example, suspensions of
Herman's minimal positive-entropy diffeomorphisms have this property; see
\cite{exa}. Thus one cannot prove Theorem~\ref{thm:main} merely by
exhibiting an atomic invariant measure supported on a periodic orbit.
\end{remark}

\section{Proof of the theorem}

\begin{proof}[Proof of Theorem~\ref{thm:main}]
We argue by contraposition. Assume that
\[
h_{\mathrm{top}}(\varphi_1)>0.
\]
By Lemma~\ref{lem:variational}, there exists an ergodic
$\varphi_t$-invariant probability measure $\mu$ with
$h_\mu(\varphi_1)>0$. By Ruelle's inequality
(Lemma~\ref{lem:ruelle}), $\mu$ has a positive Lyapunov exponent.

Therefore Lemma~\ref{lem:anosov-pesin} applies, and we obtain a
distribution
\[
u\in \cD'(M)\setminus C^\infty(M)
\]
such that
\[
\cL_X u=0.
\]
But this contradicts Lemma~\ref{lem:invariant-distributions}, which says
that under the regularity hypothesis every invariant distribution is smooth.

Hence the assumption $h_{\mathrm{top}}(\varphi_1)>0$ is impossible, and
therefore
\[
h_{\mathrm{top}}(\varphi_1)=0.
\]
This proves the theorem.
\end{proof}

\begin{thebibliography}{99}

\bibitem{KH}
A. Katok and B. Hasselblatt,
\emph{Introduction to the Modern Theory of Dynamical Systems},
Cambridge University Press, 1995.

\bibitem{Rue}
D. Ruelle,
\emph{An inequality for the entropy of differentiable maps},
Bol. Soc. Brasil. Mat. \textbf{9} (1978), 83--87.

\bibitem{Pesin}
Ya. B. Pesin,
\emph{Characteristic Lyapunov exponents and smooth ergodic theory},
Russian Math. Surveys \textbf{32} (1977), 55--114.

\bibitem{RS}
D. Ruelle and D. Sullivan,
\emph{Currents, flows and diffeomorphisms},
Topology \textbf{14} (1975), 319--329.

\bibitem{exa}
\href{https://exa.ai/library/publication/lbqpyt46khy}{exa.ai},
\emph{Topological entropies of equivalent smooth flows}.

\bibitem{ar5iv}
\href{https://ar5iv.labs.arxiv.org/html/1111.0111}{ar5iv.labs.arxiv.org},
\emph{Entropy and periodic orbits for equivalent smooth flows}.

\end{thebibliography}

\end{document}
```